\scanpage{75}
On \deleted{posera}
\[
(E,q,\omega)=(\widetilde{\mathfrak g},q,\omega)\qquad\text{i.e}\quad E=\widetilde{\mathfrak g}\tag*{(50)}
\]

\deleted{On peut aussi définir,}

\uline{NB} L'hom \deleted{\unsure{$f$ de l|fondamental}}
\[
f:\widetilde{\mathfrak g}\longrightarrow\mathfrak g\tag*{(51)}
\]
de (36) ``n'est autre'' que l'hom $\widetilde{\mathfrak g}\to\overset{\deleted{\unsure{\sim}}}{\widetilde{\mathfrak g}}^\vee$ associé à la forme bilinéaire $\varphi$, compte tenu de l'iso $\widetilde{\mathfrak g}^\vee\simeq\mathfrak g$. Ceci signifie en effet que
\[
\begin{aligned}
\langle\overset{\deleted{\unsure{w|-}}}{u},f(v)\rangle&=\varphi(u,v)\qquad u,v\in\Gamma(S,\widetilde{\mathfrak g})\\
\shortparallel\qquad&\quad\shortparallel\\
\operatorname{Tr}uv&\quad\operatorname{Tr}uv
\end{aligned}
\]

Ceci dit, les lois de crochet sur $\widetilde{\mathfrak g}$ et sur son dual $\mathfrak g$ sont celles déduites canoniquement par ailleurs des couples $(\varphi,\omega)$ d'une forme bil. sym de $\widetilde{\mathfrak g}$ et d'une base $\omega$ de $\det\widetilde{\mathfrak g}$ (c'est ici que la convention de signes pour $\omega$ importe -- on prendra
\[
\left|\overline{\omega=\widetilde X\wedge\widetilde H\wedge\widetilde Y
=-\underset{\substack{\shortparallel\\e_1}}{\widetilde X}\wedge\underset{\substack{\shortparallel\\e_2}}{\widetilde Y}\wedge\underset{\substack{\shortparallel\\e_3}}{\widetilde H}\ .}\right|\tag*{(52)}
\]
\margin{Sans garantie -- il \deleted{faut} vérifier \unsure{le texte|les textes} \ldots}
% slanted
\deleted{On}
% archive sheet number: 74
